\subsection{7.1 The Space $\mathbb{R}^m$ and Its Subsets}
\hfill 411
\subsection{7.1.2 Open and Closed Sets in $\mathbb{R}^m$}
\begin{definition} \label{def:ball} For $\delta > 0$ the set
$$
B(\mathbf{a}; \delta) = \{\mathbf{x} \in \mathbb{R}^m \mid d(\mathbf{a}, \mathbf{x}) < \delta\}
$$
is called the \emph{ball} with center $\mathbf{a} \in \mathbb{R}^m$ of radius $\delta$ or the $\delta$-neighborhood of the point
$\mathbf{a} \in \mathbb{R}^m$.
\end{definition}

\begin{definition} \label{def:open} A set $G \subset \mathbb{R}^m$ is \emph{open} in $\mathbb{R}^m$ if for every point $\mathbf{x} \in G$ there is a ball
$B(\mathbf{x}; \delta)$ such that $B(\mathbf{x}; \delta) \subset G$.
\end{definition}

\textbf{Example 1} $\mathbb{R}^m$ is an open set in $\mathbb{R}^m$.

\textbf{Example 2} The empty set $\emptyset$ contains no points at all and hence may be regarded as
satisfying Definition 2, that is, $\emptyset$ is an open set in $\mathbb{R}^m$.

\textbf{Example 3} A ball $B(\mathbf{a}; r)$ is an open set in $\mathbb{R}^m$. Indeed, if $\mathbf{x} \in B(\mathbf{a}; r)$, that is,
$d(\mathbf{a}, \mathbf{x}) < r$, then for $0 < \delta < r - d(\mathbf{a}, \mathbf{x})$, we have $B(\mathbf{x}; \delta) \subset B(\mathbf{a}; r)$, since
$$
(\mathbf{\xi} \in B(\mathbf{x}; \delta)) \implies (d(\mathbf{x}, \mathbf{\xi}) < \delta) \implies
$$
$$
\implies (d(\mathbf{a},\mathbf{\xi}) \le d(\mathbf{a}, \mathbf{x}) + d(\mathbf{x}, \mathbf{\xi}) < d(\mathbf{a}, \mathbf{x}) + r - d(\mathbf{a}, \mathbf{x}) =r).
$$

\textbf{Example 4} A set $G = \{\mathbf{x} \in \mathbb{R}^m \mid d(\mathbf{a}, \mathbf{x}) > r\}$, that is, the set of points whose dis-
tance from a fixed point $\mathbf{a} \in \mathbb{R}^m$ is larger than $r$, is open. This fact is easy to verify,
as in Example 3, using the triangle inequality for the metric.

\begin{definition} \label{def:closed} The set $F \subset \mathbb{R}^m$ is \emph{closed} in $\mathbb{R}^m$ if its complement $G = \mathbb{R}^m \setminus F$ is open
in $\mathbb{R}^m$.
\end{definition}

\textbf{Example 5} The set $\overline{B}(\mathbf{a}; r) = \{\mathbf{x} \in \mathbb{R}^m \mid d(\mathbf{a}, \mathbf{x}) \le r\}$, $r > 0$, that is, the set of points
whose distance from a fixed point $\mathbf{a} \in \mathbb{R}^m$ is at most $r$, is closed, as follows from
Definition 3 and Example 4. The set $\overline{B}(\mathbf{a}; r)$ is called the \emph{closed ball} with center $\mathbf{a}$
of radius $r$.

\begin{proposition} \label{prop:open_closed_sets} a) The union $\bigcup_{\alpha \in A} G_\alpha$ of the sets of any system $\{G_\alpha, \alpha \in A\}$ of open
sets in $\mathbb{R}^m$ is an open set in $\mathbb{R}^m$.
b) The intersection $\bigcap_{i=1}^n G_i$ of a finite number of open sets in $\mathbb{R}^m$ is an open set
in $\mathbb{R}^m$.
a') The intersection $\bigcap_{\alpha \in A} F_\alpha$ of the sets of any system $\{F_\alpha, \alpha \in A\}$ of closed sets
$F_\alpha$ in $\mathbb{R}^m$ is a closed set in $\mathbb{R}^m$.
b') The union $\bigcup_{i=1}^n F_i$ of a finite number of closed sets in $\mathbb{R}^m$ is a closed set
in $\mathbb{R}^m$.
\end{proposition}